%!TEX root = ../main.tex

\chapter*{\label{ch:introduction}Introduction}
\markboth{Introduction}{}
\addcontentsline{toc}{chapter}{Introduction} 
% NOTE: Introduction and Conclusions are special chapters: they are not
%       marked by a section number, so we need to add it manually to the
%       ToC, as done above.


Historically, theoretical and computational scientific analyses have been carried out using \textbf{explicit} mathematical models, which have grown progressively more sophisticated over time. In recent decades, however, we have witnessed a potential paradigm shift: Machine Learning (ML), and Deep Learning (DL) in particular, have enabled the \textbf{implicit} modeling of real-world phenomena, bypassing the need for many simplifying assumptions. While this has endowed Neural Networks (NNs) with unprecedented fitting power, it has also drastically compromised their interpretability.

This thesis seeks to bridge the gap between these two methodologies. Embracing the new DL paradigm should not necessitate discarding our accumulated domain knowledge; this is precisely where Scientific Machine Learning (SciML) comes into play. 
\newline
The core modeling approach of this work involves developing DL models capable of accurately emulating target systems, whether synthetic or derived from real-world data.
In this context, emulation is defined as the construction of a fast, data-driven surrogate model that faithfully reproduces the dynamics and the input-output mapping of a more complex, computationally expensive mechanistic model or physical system.

The manuscript is structured as follows: in Chapter 1 we will introduce the foundational neuroscientific ionic models that will serve as a basis for subsequent analyses. In Chapter 2 we will provide a comprehensive overview of the SciML framework. In Chapter 3, we will apply the techniques developed in the previous chapter to the emulation of neuroscientific systems using both simulated and real data. Finally, in Chapter 4 we will focus on a highly versatile mathematical framework with applications ranging from morphogenesis to disease modeling: reaction-diffusion equations, originally formalized by Alan Turing in 1952 \cite{turing1952chemical}.